\chapter{Comunicación de arquitectura: C4, UML y ADR}

\begin{objetivos}
El lector podrá seleccionar una vista según su audiencia, producir diagramas C4 coherentes, complementar con UML cuando exista una pregunta dinámica o estructural concreta y registrar decisiones mediante ADR versionados.
\end{objetivos}

\section{Una vista responde una pregunta}
Un diagrama útil declara alcance, audiencia, elementos, relaciones y leyenda. Mezclar procesos de negocio, servidores, clases y proveedores en una sola imagen genera ambigüedad. C4 organiza la estructura estática en niveles de acercamiento: contexto, contenedores, componentes y código. Los dos primeros suelen ser suficientes para discutir arquitectura de solución.

\begin{figure}[H]
\centering
\begin{tikzpicture}[every node/.style={align=center},node distance=0.75cm]
  \node[draw=AzulUD,fill=AzulClaro,rounded corners,minimum width=11.5cm,minimum height=1.1cm] (ctx) {1. Contexto: personas y sistemas externos};
  \node[draw=AzulUD,fill=white,rounded corners,below=of ctx,minimum width=9.5cm,minimum height=1.1cm] (con) {2. Contenedores: aplicaciones y almacenes de datos};
  \node[draw=Verde,fill=VerdeClaro,rounded corners,below=of con,minimum width=7.5cm,minimum height=1.1cm] (cmp) {3. Componentes: responsabilidades internas};
  \node[draw=Gris,fill=GrisClaro,rounded corners,below=of cmp,minimum width=5.5cm,minimum height=1.1cm] (cod) {4. Código: detalle selectivo};
  \draw[-{Latex},thick] (ctx) -- node[right]{acercamiento} (con);
  \draw[-{Latex},thick] (con) -- (cmp);
  \draw[-{Latex},thick] (cmp) -- (cod);
\end{tikzpicture}
\caption{Niveles de abstracción del modelo C4.}
\label{fig:c4}
\end{figure}

\section{Contexto y contenedores de LibreReserva}
En contexto aparecen Estudiante, Docente, Administrador, el sistema institucional de identidad y el servicio de correo. No se muestran frameworks ni tablas. La pregunta es quién usa el sistema y con qué sistemas se relaciona.

En contenedores se distinguen interfaz web, API, base de datos y trabajador de notificaciones. Aquí sí se indican responsabilidades, tecnologías y protocolos. ``Contenedor'' en C4 no significa necesariamente contenedor OCI; puede ser una aplicación, proceso o almacén desplegable.

\section{Cuándo emplear UML}
UML sigue siendo valioso cuando se selecciona el diagrama adecuado:
\begin{itemize}
  \item secuencia para una interacción crítica, como confirmar una reserva;
  \item estados para el ciclo solicitada--confirmada--cancelada--expirada;
  \item clases para un modelo de dominio complejo, sin convertir cada atributo en documentación permanente;
  \item despliegue para relacionar nodos, procesos y redes cuando C4 no ofrece el detalle requerido.
\end{itemize}

La consistencia importa más que la notación. Una API llamada \emph{Reservas} en el diagrama de contenedores no debería aparecer como \emph{Booking Core} en una secuencia sin explicar el cambio.

\section{ADR: decisiones que sobreviven a la reunión}
Un Architecture Decision Record registra una decisión significativa. Un formato mínimo contiene título, estado, contexto, opciones, decisión, consecuencias y fecha. Los ADR deben vivir junto al código para que cambien mediante revisión.

\begin{lstlisting}
# ADR-0001: iniciar como monolito modular

Estado: aceptado

Contexto:
Un equipo de cuatro personas debe entregar la primera version en doce semanas.

Decision:
Usar un unico despliegue con modulos Reservas, Recursos, Identidad y Notificaciones.

Consecuencias:
+ transacciones y despliegue simples;
+ menor costo operativo;
- se requieren pruebas de limites para evitar acoplamiento interno;
- la extraccion futura de un modulo exigira contratos explicitos.
\end{lstlisting}

\section{Diagramas como código}
\begin{instalacion}
Instale PlantUML y Graphviz.

Debian/Ubuntu:
\begin{lstlisting}[language=bash]
sudo apt-get install -y default-jre graphviz plantuml
\end{lstlisting}

macOS:
\begin{lstlisting}[language=bash]
brew install graphviz plantuml
\end{lstlisting}

Compruebe con \comando{plantuml -version}. También puede usar Mermaid CLI o Structurizr Lite. No instale extensiones de editor sin revisar permisos y origen.
\end{instalacion}

\begin{lstlisting}
@startuml
actor Estudiante
rectangle "LibreReserva" {
  [Interfaz web] --> [API de reservas] : HTTPS/JSON
  [API de reservas] --> database "PostgreSQL" : SQL
  [API de reservas] --> [Notificaciones] : evento
}
Estudiante --> [Interfaz web] : consulta y reserva
@enduml
\end{lstlisting}

Guarde el texto como \software{docs/contexto.puml} y genere la imagen:
\begin{lstlisting}[language=bash]
plantuml -tsvg docs/contexto.puml
\end{lstlisting}

\begin{validacion}
El SVG debe generarse sin errores. Revise automáticamente que la documentación exista y que corresponda a la fuente versionada:
\begin{lstlisting}[language=bash]
test -s docs/contexto.svg
git diff --exit-code -- docs/contexto.svg
\end{lstlisting}
\end{validacion}

\section{Ejercicios}
\begin{ejercicio}[title={Ejercicio 1. Paquete mínimo de arquitectura}]
Construya para LibreReserva un diagrama de contexto, uno de contenedores, una secuencia de confirmación y un diagrama de estados de Reserva. Cada relación debe estar etiquetada y cada tecnología debe aparecer solamente en el nivel adecuado.
\end{ejercicio}

\begin{ejercicio}[title={Ejercicio 2. Revisión cruzada}]
Intercambie los diagramas con otro equipo. Sin explicación oral, identifique usuarios, responsabilidades, protocolos, almacenes, fronteras de confianza y recorrido de una reserva. Registre todas las preguntas que el diagrama no permita responder y modifíquelo sólo cuando la pregunta pertenezca a su alcance.
\end{ejercicio}

\begin{ejercicio}[title={Ejercicio 3. ADR alternativo}]
Redacte un ADR que rechace microservicios en la primera versión. Incluya condiciones observables que justificarían reconsiderar la decisión, por ejemplo equipos independientes, ciclos de despliegue incompatibles o requisitos de aislamiento.
\end{ejercicio}
